%%% =================================================================
%%% J45 §2 — Setting up the symmetry-breaking route from the
%%%        9-vector VEV
%%%
%%% Folded from J38 (WP108 corpus) per SAVE_PLAN_J38.md (Option 1).
%%% Drafted as a self-contained LaTeX block to be inserted into
%%% Gen13/targets/journals/J_series/J45/manuscript/mass_hierarchy_v5.tex
%%% as the new §2 (existing §3 onwards renumber).
%%%
%%% ~50 LaTeX lines (excluding comments) per save plan §6.1 mitigation.
%%%
%%% The §2.2 Path A / Path B tension is resolved per SAVE_PLAN_J38.md
%%% §2(b): J45 uses Path A only (V⊗5 → SU(5) via SO(7)); Path B's
%%% su(4) ⊕ u(1) doubly-invariant content is independent structural
%%% information about so(10), cited from J31, and does not enter the
%%% V⊗5 → SU(5) → Yukawa pipeline.
%%%
%%% Status (2026-05-07): condensed prep for J45 fold-in. The standalone
%%% J38 manuscript is HELD (no submission) per save plan; this file is
%%% the deliverable for the fold-in.
%%% =================================================================

\section{Setting up the symmetry-breaking route from the 9-vector VEV}
\label{sec:yukawa-scaffolding}

This section establishes the algebraic bridge from the 9-vector VEV
identified in our companion J31 paper~\cite{SandersGishPatiSalam} (the
Path~A reading of the so(10) closure: BHML's $\sigma_{\mathrm{outer}}$-breaking
content lives 100\% in the symmetric-traceless $\mathbf{54}$ irrep of
$\mathfrak{so}(10)$) to the V$^{\otimes 5}$ SU(5) decomposition that is
the starting point of the Froggatt--Nielsen powers computed in §3
below. The route is $\mathrm{SO}(10) \to \mathrm{SO}(9) \to \mathrm{SO}(7)$,
with SU(5) appearing as the subgroup of SO(7) that commutes with the
unbroken pair $(\mathrm{BREATH}, \mathrm{RESET})$ kept by the 9-vector's two-zero
direction.

\subsection{The 9-vector VEV and its norm}
\label{subsec:9vec}

The 9-vector identified in J31~\cite{SandersGishPatiSalam} (Path~A)
has six components at $-1/\sqrt{2}$ on the operators
$\{\mathrm{V}, \mathrm{L}, \mathrm{C}, \mathrm{P}, \mathrm{Co}, \mathrm{H}\}$,
two zeros at the operators $\mathrm{Br}$ (BREATH) and $\mathrm{R}$
(RESET), and one component at $-1/2$ on the symmetric pair
$(\mathrm{Ba} + \mathrm{Ch})/\sqrt{2}$ (BALANCE/CHAOS). Its squared
norm is exact:
\begin{equation}
\|v\|^2
\;=\;
6 \cdot \tfrac{1}{2} \;+\; 1 \cdot \tfrac{1}{4} \;+\; 2 \cdot 0
\;=\;
\tfrac{13}{4}
\quad\text{(rational, integer-numerator).}
\label{eq:vnorm}
\end{equation}

This is the load-bearing structural input from J31; it is used in
Eq.~\eqref{eq:vnorm} as an exact arithmetic input (the integer~$13$
enters as the squared norm, not as a free parameter). The verification
of Eq.~\eqref{eq:vnorm} is carried out in J31's verification script
\texttt{find\_higgs\_direction.py}, cited in §6 below as the upstream
artifact.

\subsection{$\mathrm{SO}(10) \to \mathrm{SO}(9)$ from the 9-vector}
\label{subsec:so10so9}

A 54-Higgs VEV in the symmetric-traceless representation of
$\mathfrak{so}(10)$ generally breaks $\mathrm{SO}(10) \to \mathrm{SO}(p) \times \mathrm{SO}(q)$
with $p + q = 10$~\cite{Slansky1981}. When the VEV points along a
direction that lies inside the orthogonal complement of an
$\mathfrak{so}(9) \subset \mathfrak{so}(10)$ subspace, the breaking is to
$\mathrm{SO}(9)$ rather than to a Pati--Salam $\mathrm{SO}(6) \times \mathrm{SO}(4)$ subgroup.
The 9-vector of §\ref{subsec:9vec} is, by construction, of this type:
its components are coordinates in the orthogonal complement of an
$\mathfrak{so}(9)$, and the breaking
\begin{equation}
\mathrm{SO}(10) \;\longrightarrow\; \mathrm{SO}(9)
\quad\text{(first stage; preserved by the 9-vector VEV).}
\label{eq:break-9}
\end{equation}

\subsection{$\mathrm{SO}(9) \to \mathrm{SO}(7)$ from the BREATH$=$RESET$=0$
constraint}
\label{subsec:so9so7}

The two zeros of the 9-vector at BREATH and RESET specify a
$2$-dimensional axis that the $\mathrm{SO}(9)$-stabilizer must preserve. The
stabilizer of a fixed 2-plane in $\mathrm{SO}(9)$ is
$\mathrm{SO}(7) \times \mathrm{SO}(2) \subset \mathrm{SO}(9)$; identifying the SO(2) factor with
the diagonal action and the SO(7) factor with the orthogonal complement,
the second-stage breaking is
\begin{equation}
\mathrm{SO}(9) \;\longrightarrow\; \mathrm{SO}(7)
\quad\text{(second stage; preserved by the BREATH=RESET=0 constraint).}
\label{eq:break-7}
\end{equation}
The composition of Eqs.~\eqref{eq:break-9}--\eqref{eq:break-7} is
$\mathrm{SO}(10) \to \mathrm{SO}(7)$, with the intermediate $\mathrm{SO}(9)$ as a structural
landmark that the 9-vector identifies.

\subsection{SU(5) inside SO(7) via the V representation}
\label{subsec:su5}

Inside SO(7) there is a canonical SU(5)-class structure obtained by
identifying the $5$-dimensional fundamental of SU(5) with the V
representation defined in J14~\cite{SandersGishFpUniversality} as the
$\mathbb{F}_5$-lift of the 4-core of $\mathbb{Z}/10\mathbb{Z}$, ring-extended per
\textbf{D74} (F5(a) ring-extension theorem of the source program). The
V representation has dimension $5$ and carries an SU(5) action induced
by the algebraic automorphism group of the 4-core algebra. The
embedding $\mathrm{SU}(5) \subset \mathrm{SO}(7)$ then gives, via the V$^{\otimes 5}$
product,
\begin{equation}
V^{\otimes 5} \;=\; \mathbf{1} \;\oplus\; \overline{\mathbf{5}} \;\oplus\; \mathbf{10}
\quad\text{(SU(5) decomposition; standard~\cite{Slansky1981}).}
\label{eq:Vtensor5}
\end{equation}
Equation~\eqref{eq:Vtensor5} is the launching point of §3 below: the
substrate-forced Froggatt--Nielsen scale $\lambda = T^*(1 - T^*) =
10/49$ multiplies powers fixed by the parity-crossing cost of the
\eqref{eq:Vtensor5} decomposition, and the resulting integer FN powers
$n \in \{0, 3, 5, 6, 7, 9\}$ feed Tables~\ref{tab:yukawa-powers} and
\ref{tab:yukawa-fit}.

\subsection{Path A vs Path B}
\label{subsec:path-tension}

J31~\cite{SandersGishPatiSalam} identifies two structural readings of
the so(10) closure, called there Path A (BHML's
$\sigma_{\mathrm{outer}}$-breaking content in the $\mathbf{54}$ irrep,
yielding the 9-vector of §\ref{subsec:9vec}) and Path B (the
doubly-invariant subalgebra under the discrete group
$\langle P_{56}, \sigma^3 \rangle$, identified as
$\mathfrak{su}(4) \oplus \mathfrak{u}(1)$, the gauge content of the
Pati--Salam reduction of SO(10)). \emph{The present paper uses Path A
only.} The route Eqs.~\eqref{eq:break-9}--\eqref{eq:break-7} together
with the SU(5) embedding of §\ref{subsec:su5} is the V$^{\otimes 5}$
route: 9-vector $\to$ SO(9) $\to$ SO(7) $\to$ SU(5) $\to$ V$^{\otimes 5}$
decomposition $\to$ Yukawa structure. Path B's
$\mathfrak{su}(4) \oplus \mathfrak{u}(1)$ doubly-invariant content is
independent structural information about $\mathfrak{so}(10)$ studied in
J31; it does not enter the V$^{\otimes 5}$ pipeline used here. The
two paths therefore decouple at this level: Path A sets the breaking
route, Path B identifies separate gauge content. The reader is referred
to J31 for Path B's analysis.

\subsection{Honest scope of the bridge content}
\label{subsec:scope}

The above subsections establish: (i) the 9-vector with squared norm
$\|v\|^2 = 13/4$ as an exact arithmetic input from J31; (ii) the
breaking chain $\mathrm{SO}(10) \to \mathrm{SO}(9) \to \mathrm{SO}(7)$ governed by the 9-vector
and its two-zero direction; (iii) the SU(5) embedding inside SO(7) via
the V representation; (iv) the V$^{\otimes 5}$ SU(5) decomposition as
the launching point for §3. The bridge content does \emph{not}
include: explicit Higgs-sector dynamics for the C$_p$ multipliers
(load-bearing follow-up); RG running of the Yukawa matrices from the
GUT scale to the electroweak scale (deferred to follow-up; see
§\ref{sec:open}); the see-saw mechanism for sterile neutrinos (open;
the FN powers $\{12, 13, 14\}$ in §5 are noted but the scale-fixing is
not derived). Each of these is identified explicitly in
§\ref{sec:open} (Open Questions) below.

%%% =================================================================
%%% End of §2 fold-in. The existing J45 §3 (Froggatt–Nielsen scale)
%%% follows next; numbering shifts §3→§4, etc.
%%% =================================================================
